These notes are written for students who have already met mathematics in school and in earlier courses, but who may still not feel that mathematics is a language they can actively use. Many students can recognize a formula, copy a procedure, or reproduce a calculation after seeing a similar example. This is useful, but it is not yet mathematical understanding. To understand mathematics means to know what kind of object we are working with, what the symbols represent, what can be changed, what must remain fixed, what question is being asked, and why a particular operation is allowed. A formula is not only something into which numbers are inserted. A formula describes a relation, a structure, a transformation, a dependence, or a rule of change. Learning to read such objects is one of the central aims of these notes.

The course is built around three areas that are essential for later work in physics, engineering, computer science, data analysis, modelling, and many other technical disciplines: linear algebra, analytic geometry, and calculus. These areas are often taught as separate collections of definitions, formulas, and methods. Here they will be treated as parts of one larger mathematical language. Linear algebra gives us a way to speak about vectors, matrices, systems of linear equations, transformations, and structure. Analytic geometry connects algebraic equations with geometric objects such as points, lines, planes, directions, distances, and angles. Calculus gives us a language for describing change, approximation, limiting processes, accumulation, rates, and continuous models. Each of these subjects is important on its own, but their real power appears when they are used together to describe and analyse a situation.

The purpose of these notes is therefore not to present mathematics as a dry catalogue of definitions and theorems. Definitions and theorems are necessary, but they are not sufficient. A definition introduces an object, but the student must also learn how to look at that object, how to test it, how to transform it, how to compare it with other objects, and how to ask useful questions about it. A theorem gives a reliable statement, but the student must also understand what the assumptions mean, why they are needed, and how the theorem changes the way we can reason. An example is not merely a pattern to imitate mechanically. A good example should show how a concept behaves, what choices are being made, and how a calculation expresses an idea.

For this reason, the text will often move slowly through ideas that may at first seem elementary. We will not assume that a student can read a formula just because the formula has been written on the page. When we write something as simple as
\[
y = x^2,
\]
we should be able to say what is variable, what is being assigned, what happens when a value of \(x\) is chosen, how the value of \(y\) is produced, and how the same relation can be seen as a graph, a rule, a table of values, or a model of dependence. When we solve an equation, we should know whether we are finding a number, a vector, a point, a function, or a set of possible objects. When we multiply a matrix by a vector, we should not see only an array of arithmetic operations; we should also see a transformation of information. When we take a derivative, we should not see only a formal rule; we should also see a rate of change and a local approximation. When we take an integral, we should not see only an antiderivative; we should also see accumulation, total effect, area, displacement, or another quantity built from many small contributions.

The notes will repeatedly return to the same basic habits of thought. First we identify the object. Then we ask what information is given and what is unknown. We decide which notation is useful. We translate the question into mathematical language. We perform operations carefully, but we also ask what those operations mean. After obtaining a result, we interpret it: we check whether the type of the answer is correct, whether the units or dimensions make sense when they are present, whether the result is reasonable, and whether special or limiting cases agree with intuition. These habits are intentionally repeated. Repetition is not a weakness of the exposition. It is a way of building fluency. A student who sees the same style of reasoning in vectors, functions, equations, derivatives, and integrals should gradually begin to recognize mathematics as a coherent language rather than as a set of disconnected tricks.

This approach is especially important because these notes prepare the ground for physics. In physics, mathematics is not used as decoration. It is the language in which models are built. A physical law written as an equation does not merely tell us which numbers to substitute. It tells us what quantities are related, what assumptions are being made, what has been idealized, and what kind of prediction can be made. Students who cannot read functions, vectors, equations, and derivatives as meaningful mathematical objects will later have difficulty understanding motion, force, energy, fields, waves, and other physical ideas. The mathematical part of the course must therefore develop not only computational skill, but also the ability to interpret mathematical expressions as statements about structure and change.

This does not mean that the course avoids calculation. On the contrary, calculation is necessary. But calculation should not be separated from meaning. Algebraic manipulation is valuable when the student understands what is being preserved and what is being transformed. Solving a system of equations is valuable when the student knows what the unknowns represent and what the solution means. Drawing a line or a plane is valuable when the equation is connected with direction, normal vectors, and distance. Differentiating a function is valuable when the derivative is connected with local behaviour, monotonicity, approximation, and rates of change. Integrating a function is valuable when the integral is connected with accumulation and total quantity. The goal is not to remove techniques, but to embed them in reasoning.

The reader should therefore study these notes actively. Do not only ask: Which formula should I use? Ask instead: What kind of object is this? What is known? What is unknown? What does this symbol mean here? What changes if this variable changes? What remains fixed? What question is being asked? What operation is allowed, and why? What does the result say? Could the answer have been expected from a graph, a simple example, or a limiting case? Such questions may seem slow at first, but they are the beginning of real mathematical independence.

The aim of the course is to help the student move from passive recognition to active use. A student should become able not only to follow a solved example, but also to read a new mathematical expression, identify its structure, ask meaningful questions, choose a method, carry out a calculation, and explain the result. This is the level of mathematical literacy needed for later physics and for technical work more generally. Mathematics in these notes is not presented as an obstacle to be survived. It is presented as a language for thinking clearly about relations, structures, models, and change.
